% Options for packages loaded elsewhere
\PassOptionsToPackage{unicode}{hyperref}
\PassOptionsToPackage{hyphens}{url}
\PassOptionsToPackage{space}{xeCJK}
%
\documentclass[
]{ctexart}
\usepackage{amsmath,amssymb}

\usepackage{iftex}
\ifPDFTeX
  \usepackage[T1]{fontenc}
  \usepackage[utf8]{inputenc}
  \usepackage{textcomp} % provide euro and other symbols
\else % if luatex or xetex
  \usepackage{unicode-math}
  \defaultfontfeatures{Scale=MatchLowercase}
  \defaultfontfeatures[\rmfamily]{Ligatures=TeX,Scale=1}
  \ifXeTeX
    \usepackage{xeCJK}
    \setCJKmainfont[]{FangSong}
  \fi
  \ifLuaTeX
    \usepackage[]{luatexja-fontspec}
    \setmainjfont[]{FangSong}
  \fi
\fi
% Use upquote if available, for straight quotes in verbatim environments
\IfFileExists{upquote.sty}{\usepackage{upquote}}{}
\IfFileExists{microtype.sty}{% use microtype if available
  \usepackage[]{microtype}
  \UseMicrotypeSet[protrusion]{basicmath} % disable protrusion for tt fonts
}{}
\makeatletter
\@ifundefined{KOMAClassName}{% if non-KOMA class
  \IfFileExists{parskip.sty}{%
    \usepackage{parskip}
  }{% else
    \setlength{\parindent}{0pt}
    \setlength{\parskip}{6pt plus 2pt minus 1pt}}
}{% if KOMA class
  \KOMAoptions{parskip=half}}
\makeatother
\usepackage{xcolor}
\IfFileExists{xurl.sty}{\usepackage{xurl}}{} % add URL line breaks if available
\IfFileExists{bookmark.sty}{\usepackage{bookmark}}{\usepackage{hyperref}}
\hypersetup{
  hidelinks}
\urlstyle{same} % disable monospaced font for URLs
\usepackage[margin=2.1cm,headheight=15pt]{geometry}
\usepackage{color}
\usepackage{fancyvrb}
\newcommand{\VerbBar}{|}
\newcommand{\VERB}{\Verb[commandchars=\\\{\}]}
\DefineVerbatimEnvironment{Highlighting}{Verbatim}{commandchars=\\\{\}}
% Add ',fontsize=\small' for more characters per line
\newenvironment{Shaded}{}{}
\newcommand{\AlertTok}[1]{\textcolor[rgb]{1.00,0.00,0.00}{\textbf{#1}}}
\newcommand{\AnnotationTok}[1]{\textcolor[rgb]{0.38,0.63,0.69}{\textbf{\textit{#1}}}}
\newcommand{\AttributeTok}[1]{\textcolor[rgb]{0.49,0.56,0.16}{#1}}
\newcommand{\BaseNTok}[1]{\textcolor[rgb]{0.25,0.63,0.44}{#1}}
\newcommand{\BuiltInTok}[1]{#1}
\newcommand{\CharTok}[1]{\textcolor[rgb]{0.25,0.44,0.63}{#1}}
\newcommand{\CommentTok}[1]{\textcolor[rgb]{0.38,0.63,0.69}{\textit{#1}}}
\newcommand{\CommentVarTok}[1]{\textcolor[rgb]{0.38,0.63,0.69}{\textbf{\textit{#1}}}}
\newcommand{\ConstantTok}[1]{\textcolor[rgb]{0.53,0.00,0.00}{#1}}
\newcommand{\ControlFlowTok}[1]{\textcolor[rgb]{0.00,0.44,0.13}{\textbf{#1}}}
\newcommand{\DataTypeTok}[1]{\textcolor[rgb]{0.56,0.13,0.00}{#1}}
\newcommand{\DecValTok}[1]{\textcolor[rgb]{0.25,0.63,0.44}{#1}}
\newcommand{\DocumentationTok}[1]{\textcolor[rgb]{0.73,0.13,0.13}{\textit{#1}}}
\newcommand{\ErrorTok}[1]{\textcolor[rgb]{1.00,0.00,0.00}{\textbf{#1}}}
\newcommand{\ExtensionTok}[1]{#1}
\newcommand{\FloatTok}[1]{\textcolor[rgb]{0.25,0.63,0.44}{#1}}
\newcommand{\FunctionTok}[1]{\textcolor[rgb]{0.02,0.16,0.49}{#1}}
\newcommand{\ImportTok}[1]{#1}
\newcommand{\InformationTok}[1]{\textcolor[rgb]{0.38,0.63,0.69}{\textbf{\textit{#1}}}}
\newcommand{\KeywordTok}[1]{\textcolor[rgb]{0.00,0.44,0.13}{\textbf{#1}}}
\newcommand{\NormalTok}[1]{#1}
\newcommand{\OperatorTok}[1]{\textcolor[rgb]{0.40,0.40,0.40}{#1}}
\newcommand{\OtherTok}[1]{\textcolor[rgb]{0.00,0.44,0.13}{#1}}
\newcommand{\PreprocessorTok}[1]{\textcolor[rgb]{0.74,0.48,0.00}{#1}}
\newcommand{\RegionMarkerTok}[1]{#1}
\newcommand{\SpecialCharTok}[1]{\textcolor[rgb]{0.25,0.44,0.63}{#1}}
\newcommand{\SpecialStringTok}[1]{\textcolor[rgb]{0.73,0.40,0.53}{#1}}
\newcommand{\StringTok}[1]{\textcolor[rgb]{0.25,0.44,0.63}{#1}}
\newcommand{\VariableTok}[1]{\textcolor[rgb]{0.10,0.09,0.49}{#1}}
\newcommand{\VerbatimStringTok}[1]{\textcolor[rgb]{0.25,0.44,0.63}{#1}}
\newcommand{\WarningTok}[1]{\textcolor[rgb]{0.38,0.63,0.69}{\textbf{\textit{#1}}}}
\usepackage{longtable,booktabs,array}
\usepackage{calc} % for calculating minipage widths
% Correct order of tables after \paragraph or \subparagraph
\usepackage{etoolbox}
\makeatletter
\patchcmd\longtable{\par}{\if@noskipsec\mbox{}\fi\par}{}{}
\makeatother
% Allow footnotes in longtable head/foot
\IfFileExists{footnotehyper.sty}{\usepackage{footnotehyper}}{\usepackage{footnote}}
\makesavenoteenv{longtable}
\setlength{\emergencystretch}{3em} % prevent overfull lines
\providecommand{\tightlist}{%
  \setlength{\itemsep}{0pt}\setlength{\parskip}{0pt}}
\setcounter{secnumdepth}{5}
\usepackage{xcolor}
\usepackage{fontspec}
\usepackage{graphicx}
\usepackage{pdflscape}
\usepackage{float}
\usepackage{booktabs}
\usepackage{longtable}
\usepackage{array}
\usepackage{tabularx}
\usepackage{etoolbox}
\usepackage{fancyhdr}
\usepackage{titlesec}
\usepackage{caption}
\usepackage{enumitem}
\usepackage{fvextra}
\usepackage{setspace}

\definecolor{cfInk}{HTML}{1D1B18}
\definecolor{cfText}{HTML}{2E2A24}
\definecolor{cfMuted}{HTML}{68625A}
\definecolor{cfGold}{HTML}{B88A35}
\definecolor{cfLine}{HTML}{E4D8C6}
\definecolor{cfSoft}{HTML}{FBF7EF}
\IfFontExistsTF{FangSong}{\setmainfont{FangSong}}{\IfFontExistsTF{STFangsong}{\setmainfont{STFangsong}}{}}
\IfFontExistsTF{Microsoft YaHei}{\setsansfont{Microsoft YaHei}}{}
\IfFontExistsTF{Consolas}{\setmonofont{Consolas}[Scale=0.88]}{}
\IfFontExistsTF{FangSong}{\setCJKmainfont{FangSong}}{\IfFontExistsTF{STFangsong}{\setCJKmainfont{STFangsong}}{}}
\IfFontExistsTF{Microsoft YaHei}{\setCJKsansfont{Microsoft YaHei}}{}
\IfFontExistsTF{SimSun}{\setCJKmonofont{SimSun}}{}

\linespread{1.08}
\setlength{\parskip}{0.38em}
\setlength{\parindent}{2em}
\setlength{\emergencystretch}{4em}
\sloppy

\pagestyle{fancy}
\fancyhf{}
\fancyhead[L]{\small\sffamily CinemaFlow 电影库管理系统汇报}
\fancyhead[R]{\small\sffamily\leftmark}
\fancyfoot[C]{\small\sffamily\thepage}
\renewcommand{\headrulewidth}{0.4pt}
\renewcommand{\footrulewidth}{0pt}
\setlength{\headheight}{15pt}

\titleformat{\section}
  {\sffamily\bfseries\Large\color{cfInk}}
  {}
  {0pt}
  {}
  [\vspace{0.25em}{\color{cfGold}\titlerule[0.8pt]}]
\titleformat{\subsection}
  {\sffamily\bfseries\large\color{cfInk}}
  {}
  {0pt}
  {}
\titleformat{\subsubsection}
  {\sffamily\bfseries\normalsize\color{cfInk}}
  {}
  {0pt}
  {}
\titlespacing*{\section}{0pt}{1.1em}{0.8em}
\titlespacing*{\subsection}{0pt}{0.9em}{0.45em}
\titlespacing*{\subsubsection}{0pt}{0.7em}{0.35em}

\captionsetup{
  font=small,
  labelfont={bf,color=cfGold},
  textfont={color=cfMuted},
  skip=5pt
}

\setlist[itemize]{leftmargin=2em,itemsep=0.2em,topsep=0.25em,label={\color{cfGold}\raisebox{0.15ex}{\rule{0.36em}{0.36em}}}}
\setlist[enumerate]{leftmargin=2em,itemsep=0.2em,topsep=0.25em}

\AtBeginEnvironment{longtable}{\small\setlength{\tabcolsep}{4pt}\renewcommand{\arraystretch}{1.22}}
\AtBeginEnvironment{tabular}{\small\setlength{\tabcolsep}{4pt}\renewcommand{\arraystretch}{1.18}}
\AtBeginEnvironment{Shaded}{\small}
\fvset{breaklines=true,breakanywhere=true,fontsize=\small}

\renewcommand{\arraystretch}{1.16}
\setlength{\LTpre}{0.5em}
\setlength{\LTpost}{0.7em}

\AtBeginDocument{
  \color{cfText}
  \hypersetup{
    colorlinks=true,
    linkcolor=cfGold,
    urlcolor=cfGold,
    citecolor=cfGold,
    pdftitle={CinemaFlow 电影库管理系统汇报},
    pdfauthor={CinemaFlow 项目组},
    pdfcreator={CinemaFlow 项目组}
  }
}
\ifLuaTeX
  \usepackage{selnolig}  % disable illegal ligatures
\fi

\author{}
\date{}

\begin{document}

\begin{titlepage}
\thispagestyle{empty}
\sffamily
\color{cfText}
\vspace*{1.7cm}
\begin{center}
{\color{cfGold}\rule{0.2\linewidth}{1.2pt}}\par
\vspace{1.05cm}
{\fontsize{33}{41}\selectfont\bfseries CinemaFlow\par}
\vspace{0.35cm}
{\fontsize{24}{32}\selectfont\bfseries 电影库管理系统汇报\par}
\vspace{0.8cm}
{\Large 前后端分离架构、数据库设计与功能说明\par}
\end{center}
\vspace{1.35cm}
\begingroup
\setlength{\tabcolsep}{0pt}
\renewcommand{\arraystretch}{1.32}
\begin{center}
\begin{tabular}{>{\color{cfGold}\bfseries}r@{\hspace{0.9cm}}l}
系统类型 & 私人电影资料馆与观影管理系统\\
技术栈 & Angular 17 Standalone Components、Flask REST API、关系型数据库设计\\
文档范围 & 系统业务、架构改造、数据库、API、前端设计、核心代码与页面说明\\
版本日期 & 2026 年 6 月\\
\end{tabular}
\end{center}
\endgroup
\vfill
\noindent{\color{cfGold}\rule{\linewidth}{0.7pt}}\par
\vspace{0.35cm}
\begin{center}
{\large 课程项目汇报文档\par}
\end{center}
\end{titlepage}
\color{cfText}
\clearpage

{
\setcounter{tocdepth}{2}
\tableofcontents
}
\hypertarget{ux7cfbux7edfux7b80ux4ecb}{%
\section{系统简介}\label{ux7cfbux7edfux7b80ux4ecb}}

\hypertarget{ux4e1aux52a1ux80ccux666f}{%
\subsection{业务背景}\label{ux4e1aux52a1ux80ccux666f}}

CinemaFlow
面向个人影迷、电影资料整理者和课程项目展示场景。普通表格或笔记软件可以记录片名和评分，但当电影数量增多后，用户会自然遇到几个问题：电影资料分散、导演与作品关系不清晰、观影计划和观影日志没有闭环、收藏与偏好难以统计、媒体图片
URL 质量无法统一检查。

CinemaFlow
将电影资料、导演档案、观影计划、观影日志、影评、推荐、偏好画像和片库审计整合到一个系统里。系统不只是展示电影卡片，而是围绕``电影资料长期维护''和``观影行为沉淀''建立完整业务链路。

\hypertarget{ux4e3bux8981ux7528ux6237}{%
\subsection{主要用户}\label{ux4e3bux8981ux7528ux6237}}

\begin{longtable}[]{@{}lll@{}}
\toprule
用户角色 & 关注点 & 典型操作 \\
\midrule
\endhead
普通影迷 & 快速找到想看的电影，记录已看、收藏和评分 &
浏览电影库、收藏电影、写观影日志、查看推荐 \\
资料维护者 & 保证电影资料、导演资料、海报和背景图完整 & 添加电影、维护
URL、查看片库审计 \\
电影内容整理者 & 通过导演、类型、年代、情绪组织片单 &
使用导演图谱、情绪图谱、氛围策展、连看规划 \\
系统评审人员 & 了解系统功能完整度、架构边界和数据关系 & 查看架构图、ER
图、API 表、页面截图 \\
开发维护者 & 扩展接口、数据库和前端页面 &
阅读服务层、状态流、路由和数据库设计 \\
\bottomrule
\end{longtable}

\hypertarget{ux7cfbux7edfux89e3ux51b3ux7684ux95eeux9898}{%
\subsection{系统解决的问题}\label{ux7cfbux7edfux89e3ux51b3ux7684ux95eeux9898}}

\begin{longtable}[]{@{}ll@{}}
\toprule
问题 & CinemaFlow 的解决方式 \\
\midrule
\endhead
电影数量增加后难以查找 &
提供电影库搜索、类型路由、排序、分页和命令面板 \\
电影与导演资料割裂 & 建立电影、导演、代表作、导演详情和导演图谱 \\
``想看''和``已看''没有闭环 &
提供待看片单、观影日志、观影日历和偏好画像 \\
收藏、评分、影评难以统计 & 提供收藏中心、影评墙、用户评分和推荐算法 \\
媒体 URL 质量参差不齐 & 提供真实 URL 规范化、加载兜底和片库审计 \\
页面多、状态复杂 & 使用 Angular Router、服务层封装和 RxJS
状态流统一管理 \\
后续扩展数据库困难 &
按关系型模型设计电影、人员、日志、计划、审计等实体 \\
\bottomrule
\end{longtable}

\begin{center}\rule{0.5\linewidth}{0.5pt}\end{center}

\hypertarget{ux7cfbux7edfux529fux80fdux603bux89c8}{%
\section{系统功能总览}\label{ux7cfbux7edfux529fux80fdux603bux89c8}}

\hypertarget{ux6838ux5fc3ux4e1aux52a1ux94feux8def}{%
\subsection{核心业务链路}\label{ux6838ux5fc3ux4e1aux52a1ux94feux8def}}

\begin{center}
\centering
\includegraphics[width=0.96\linewidth,height=0.56\textheight,keepaspectratio]{output/pdf/diagrams/report-diagram-01.png}
\captionof{figure}{核心业务链路}
\end{center}

这条链路体现系统的核心价值：电影资料不是静态列表，而是会随着用户行为不断沉淀。用户从``发现一部电影''开始，可以进入详情、收藏、计划观看、完成观看、写下日志和影评，最后形成偏好画像，并反过来影响推荐和智能选片。

\hypertarget{ux529fux80fdux6a21ux5757ux5212ux5206}{%
\subsection{功能模块划分}\label{ux529fux80fdux6a21ux5757ux5212ux5206}}

\begin{longtable}[]{@{}lll@{}}
\toprule
模块 & 页面/组件 & 业务价值 \\
\midrule
\endhead
片库总览 & 仪表盘、最近浏览、最近添加 &
快速了解片库规模、评分、收藏和浏览动态 \\
电影管理 & 电影库、分类浏览、电影详情、演员表、添加电影 &
完成电影资料的浏览、检索、录入和详情展示 \\
导演资料 & 导演库、导演详情、导演图谱 &
以导演为维度组织作品和作者风格 \\
观影行为 & 待看片单、观影日志、观影日历 &
管理从计划到观看完成的行为记录 \\
内容表达 & 影评墙、收藏中心 & 保存主观评价、评论内容和高价值收藏 \\
推荐分析 & 推荐页、随机选片、智能选片、偏好画像 &
帮助用户做下一部电影选择 \\
高级策展 & 时间线、电影对比、情绪图谱、连看规划、氛围策展 &
从年代、情绪、主题和时长角度重组片库 \\
系统维护 & 片库审计、关于页、数据管理、命令面板 &
提供资料质量检查、运行状态、导入导出和快速跳转 \\
\bottomrule
\end{longtable}

\hypertarget{ux6570ux636eux5bf9ux8c61}{%
\subsection{数据对象}\label{ux6570ux636eux5bf9ux8c61}}

系统围绕以下业务对象展开：

\begin{itemize}
\tightlist
\item
  电影：片名、导演、类型、评分、时长、语言、简介、海报、背景图、预告片。
\item
  导演：姓名、国籍、出生年份、头像、活跃年代、风格、获奖、代表作。
\item
  用户电影状态：是否已看、是否收藏、个人评分、个人笔记、观看次数。
\item
  影评：作者、评分、标题、正文、点赞数。
\item
  待看片单：优先级、计划日期、观看场景、状态、备注。
\item
  观影日志：观看日期、地点、陪伴对象、会话评分、情绪标签、笔记。
\item
  媒体资源：电影海报、背景图、导演头像、提供方、检查时间。
\item
  最近浏览：用户进入详情页后的访问轨迹。
\item
  智能选片预设：时长、评分、类型、语言、是否包含已看、是否偏好收藏。
\end{itemize}

\begin{center}\rule{0.5\linewidth}{0.5pt}\end{center}

\hypertarget{ux67b6ux6784ux8bbeux8ba1ux8bf4ux660e}{%
\section{架构设计说明}\label{ux67b6ux6784ux8bbeux8ba1ux8bf4ux660e}}

\hypertarget{flask-ux5355ux4f53-bs-ux67b6ux6784ux7684ux7279ux70b9}{%
\subsection{Flask 单体 B/S
架构的特点}\label{flask-ux5355ux4f53-bs-ux67b6ux6784ux7684ux7279ux70b9}}

传统 Flask 单体 B/S 架构通常由 Flask
同时承担页面渲染、接口、静态资源、表单提交、业务逻辑和数据访问。

\begin{center}
\centering
\includegraphics[width=0.96\linewidth,height=0.56\textheight,keepaspectratio]{output/pdf/diagrams/report-diagram-02.png}
\captionof{figure}{Flask 单体 B/S 架构}
\end{center}

这种架构适合页面数量少、交互简单的系统。对于 CinemaFlow
这种页面数量多、局部状态多、搜索筛选频繁、需要命令面板和多种分析页面的系统，单体模板会让前端交互和后端数据逻辑紧密耦合。页面变化会影响后端模板，接口也不容易复用给其他客户端。

\hypertarget{ux524dux540eux7aefux5206ux79bbux67b6ux6784}{%
\subsection{前后端分离架构}\label{ux524dux540eux7aefux5206ux79bbux67b6ux6784}}

CinemaFlow 采用 Angular SPA + Flask REST API
的前后端分离结构。前端负责路由、组件、状态流和交互体验；后端负责 REST
API、数据校验、跨域策略和持久化；数据库负责保存电影资料和用户行为。

\begin{center}
\centering
\includegraphics[width=0.96\linewidth,height=0.56\textheight,keepaspectratio]{output/pdf/diagrams/report-diagram-03.png}
\captionof{figure}{前后端分离架构}
\end{center}

\hypertarget{ux67b6ux6784ux804cux8d23ux5bf9ux6bd4}{%
\subsection{架构职责对比}\label{ux67b6ux6784ux804cux8d23ux5bf9ux6bd4}}

\begin{longtable}[]{@{}lll@{}}
\toprule
维度 & Flask 单体 B/S & CinemaFlow 前后端分离 \\
\midrule
\endhead
页面渲染 & Flask 模板生成 HTML & Angular 组件在浏览器渲染 \\
页面路由 & 后端 route 控制页面跳转 & Angular Router 控制 SPA
页面和子路由 \\
数据交换 & 模板上下文或表单提交 & REST JSON \\
状态管理 & 依赖页面刷新和服务端上下文 &
RxJS、BehaviorSubject、combineLatest、LocalStorage \\
前端部署 & 依赖后端应用一起发布 & 前端构建产物可独立托管 \\
后端职责 & 页面、静态资源、接口混合 & API、校验、持久化、CORS \\
数据库演进 & 常和模板逻辑缠在一起 & API
合约稳定，数据库可在后端内部演进 \\
扩展客户端 & 复用困难 & REST API 可提供给移动端、小程序或其他客户端 \\
\bottomrule
\end{longtable}

\hypertarget{ux7cfbux7edfux5206ux5c42}{%
\subsection{系统分层}\label{ux7cfbux7edfux5206ux5c42}}

系统按职责拆成七个层次，各层之间通过清晰的数据对象和服务接口协作，避免页面组件直接处理持久化细节。

\begin{itemize}
\tightlist
\item
  表现层：对应 \texttt{pages/} 和
  \texttt{components/}，负责页面布局、用户交互、表格、卡片、图谱和表单展示。
\item
  路由层：对应 \texttt{app.routes.ts} 和
  \texttt{auth.guard.ts}，负责页面路由、详情页子路由和添加页权限守卫。
\item
  状态层：对应
  \texttt{MovieStateService}，聚合电影、最近浏览、影评、日志、计划等数据，形成页面
  ViewModel。
\item
  业务服务层：对应
  \texttt{MovieService}、\texttt{DirectorService}、\texttt{WatchLogService}
  等服务，负责实体管理、查询封装、API 同步和本地降级。
\item
  配置层：对应 \texttt{api.config.ts} 和
  \texttt{runtime-config.js}，负责解耦构建产物和 API 根路径。
\item
  后端 API 层：对应 Flask Blueprints，提供电影和导演 REST
  API，并负责请求校验、CORS 和持久化入口。
\item
  数据层：对应 JSON 演示持久层和关系型
  schema，保存电影、导演、用户行为、媒体资源和审计数据。
\end{itemize}

\begin{center}\rule{0.5\linewidth}{0.5pt}\end{center}

\hypertarget{ux6570ux636eux5e93ux8bbeux8ba1}{%
\section{数据库设计}\label{ux6570ux636eux5e93ux8bbeux8ba1}}

\hypertarget{ux8bbeux8ba1ux539fux5219}{%
\subsection{设计原则}\label{ux8bbeux8ba1ux539fux5219}}

CinemaFlow 的数据库采用关系型设计，核心原则如下：

\begin{itemize}
\tightlist
\item
  电影和导演属于主数据，字段稳定、复用频率高。
\item
  类型、演员、媒体资源拆成独立表，避免在电影表中堆叠多值字段。
\item
  收藏、已看、个人评分、笔记属于用户和电影之间的状态，因此放入
  \texttt{user\_movie\_states}。
\item
  影评、待看片单、观影日志属于用户围绕电影产生的业务事件，应独立建表。
\item
  情绪标签是观影日志的多值属性，使用 \texttt{watch\_log\_moods} 拆表。
\item
  最近浏览和审计日志属于行为轨迹，支持统计和问题追溯。
\item
  媒体资源表用于记录海报、背景图、头像等 URL，便于审计和替换。
\end{itemize}

\clearpage
\begin{landscape}
\phantomsection
\addcontentsline{toc}{subsection}{完整 ER 图}
{\sffamily\bfseries\large\color{cfInk}完整 ER 图\par}
\vspace{0.35cm}
\begin{center}
\centering
\includegraphics[width=0.98\linewidth,height=0.86\textheight,keepaspectratio]{output/pdf/diagrams/report-diagram-04.png}
\captionof{figure}{完整 ER 图}
\end{center}
\end{landscape}
\clearpage

完整 ER
图展示所有实体之间的全局关系。由于系统实体较多，为了让评审时更容易看清业务边界，下面再按``核心资料''``用户行为''``媒体与审计''拆成三张局部
ER 图。

\hypertarget{ux6838ux5fc3ux8d44ux6599-er-ux56fe}{%
\subsection{核心资料 ER 图}\label{ux6838ux5fc3ux8d44ux6599-er-ux56fe}}

\begin{center}
\centering
\includegraphics[width=0.96\linewidth,height=0.56\textheight,keepaspectratio]{output/pdf/diagrams/report-diagram-05.png}
\captionof{figure}{核心资料 ER 图}
\end{center}

核心资料区解决``电影是什么、由谁导演、属于哪些类型、有哪些演员''的问题。\texttt{movies}
是核心表，\texttt{directors} 是电影的一对多上游实体，\texttt{genres} 与
\texttt{people}
通过关系表和电影连接。这样设计可以避免电影表里出现长字符串数组，方便后续按导演、类型、演员做查询和统计。

\hypertarget{ux7528ux6237ux884cux4e3a-er-ux56fe}{%
\subsection{用户行为 ER 图}\label{ux7528ux6237ux884cux4e3a-er-ux56fe}}

\begin{center}
\centering
\includegraphics[width=0.96\linewidth,height=0.56\textheight,keepaspectratio]{output/pdf/diagrams/report-diagram-06.png}
\captionof{figure}{用户行为 ER 图}
\end{center}

用户行为区解决``用户如何使用电影资料''的问题。收藏、已看、个人评分和笔记不直接放在
\texttt{movies} 表里，而是放在
\texttt{user\_movie\_states}，这样同一部电影可以被不同用户以不同方式标记。影评、待看片单和观影日志是三类独立事件，分别对应评论、计划和实际观看。

\hypertarget{ux5a92ux4f53ux4e0eux5ba1ux8ba1-er-ux56fe}{%
\subsection{媒体与审计 ER
图}\label{ux5a92ux4f53ux4e0eux5ba1ux8ba1-er-ux56fe}}

\begin{center}
\centering
\includegraphics[width=0.96\linewidth,height=0.56\textheight,keepaspectratio]{output/pdf/diagrams/report-diagram-07.png}
\captionof{figure}{媒体与审计 ER 图}
\end{center}

媒体与审计区解决``资料是否可靠、操作是否可追溯''的问题。\texttt{media\_assets}
通过 \texttt{entity\_type\ +\ entity\_id}
关联电影、导演、人物等多类实体，适合统一管理海报、背景图、头像和检查时间。\texttt{audit\_logs}
记录关键写操作，保存修改前后数据，方便后续排查误删、误改和资料质量问题。

\hypertarget{ux4e3bux8981ux8868ux7ed3ux6784ux8bf4ux660e}{%
\subsection{主要表结构说明}\label{ux4e3bux8981ux8868ux7ed3ux6784ux8bf4ux660e}}

\begin{longtable}[]{@{}
  >{\raggedright\arraybackslash}p{(\columnwidth - 6\tabcolsep) * \real{0.25}}
  >{\raggedright\arraybackslash}p{(\columnwidth - 6\tabcolsep) * \real{0.25}}
  >{\raggedright\arraybackslash}p{(\columnwidth - 6\tabcolsep) * \real{0.25}}
  >{\raggedright\arraybackslash}p{(\columnwidth - 6\tabcolsep) * \real{0.25}}@{}}
\toprule
表 & 业务含义 & 关键字段 & 关系说明 \\
\midrule
\endhead
\texttt{users} & 系统用户 & \texttt{username}, \texttt{password\_hash},
\texttt{role} & 与评分、收藏、影评、日志、计划关联 \\
\texttt{directors} & 导演主档 & \texttt{name}, \texttt{portrait\_url},
\texttt{signature\_style} & 一位导演可执导多部电影 \\
\texttt{movies} & 电影主档 & \texttt{title}, \texttt{director\_id},
\texttt{rating}, \texttt{poster\_url} & 系统核心实体 \\
\texttt{genres} & 类型字典 & \texttt{name}, \texttt{sort\_order} &
规范类型名称和排序 \\
\texttt{movie\_genres} & 电影和类型关系 & \texttt{movie\_id},
\texttt{genre\_id} & 支持一部电影多个类型 \\
\texttt{people} & 演职人员主档 & \texttt{name}, \texttt{portrait\_url},
\texttt{bio} & 与电影通过 \texttt{movie\_people} 关联 \\
\texttt{movie\_people} & 电影演职人员关系 & \texttt{role\_type},
\texttt{character\_name}, \texttt{sort\_order} &
支持演员、编剧、摄影等角色扩展 \\
\texttt{media\_assets} & 媒体资源 & \texttt{entity\_type},
\texttt{asset\_type}, \texttt{url} & 保存海报、背景图、头像等 URL \\
\texttt{user\_movie\_states} & 用户对电影的状态 & \texttt{is\_watched},
\texttt{is\_favorite}, \texttt{user\_rating} &
解决多用户下收藏和评分隔离 \\
\texttt{reviews} & 影评 & \texttt{author}, \texttt{rating},
\texttt{content}, \texttt{likes} & 影评属于电影，可绑定用户 \\
\texttt{watch\_plans} & 待看片单 & \texttt{priority},
\texttt{planned\_date}, \texttt{status} & 记录计划观看行为 \\
\texttt{watch\_logs} & 观影日志 & \texttt{watched\_at},
\texttt{location}, \texttt{session\_rating} & 记录真实观看事件 \\
\texttt{watch\_log\_moods} & 日志情绪标签 & \texttt{watch\_log\_id},
\texttt{mood\_tag} & 一条日志可有多个情绪 \\
\texttt{recent\_history} & 最近浏览 & \texttt{viewed\_at},
\texttt{source\_route} & 支持最近访问轨迹 \\
\texttt{smart\_pick\_presets} & 智能选片预设 & \texttt{max\_duration},
\texttt{min\_rating}, \texttt{genres\_json} & 保存用户常用筛选条件 \\
\texttt{audit\_logs} & 操作审计 & \texttt{action},
\texttt{before\_json}, \texttt{after\_json} & 记录关键修改，便于追溯 \\
\bottomrule
\end{longtable}

\hypertarget{ux5173ux8054ux5173ux7cfbux8bf4ux660e}{%
\subsection{关联关系说明}\label{ux5173ux8054ux5173ux7cfbux8bf4ux660e}}

\begin{itemize}
\tightlist
\item
  \texttt{directors} 到 \texttt{movies}
  是一对多：一位导演可以对应多部电影。
\item
  \texttt{movies} 和 \texttt{genres} 是多对多：通过
  \texttt{movie\_genres} 支持一部电影同时属于科幻、动作、剧情等类型。
\item
  \texttt{movies} 和 \texttt{people} 是多对多：通过
  \texttt{movie\_people} 记录演员、角色名和排序。
\item
  \texttt{users} 和 \texttt{movies} 是多对多状态关系：通过
  \texttt{user\_movie\_states} 保存用户自己的收藏、已看、评分和笔记。
\item
  \texttt{watch\_logs} 和 \texttt{watch\_log\_moods}
  是一对多：一条观影日志可被标注为治愈、热血、失眠夜等多个情绪。
\item
  \texttt{media\_assets} 使用 \texttt{entity\_type\ +\ entity\_id}
  支持电影、导演、人物等多种实体的媒体资源管理。
\item
  \texttt{audit\_logs} 使用 \texttt{entity\_type\ +\ entity\_id}
  记录不同实体的创建、更新和删除行为。
\end{itemize}

\hypertarget{ux7d22ux5f15ux8bbeux8ba1}{%
\subsection{索引设计}\label{ux7d22ux5f15ux8bbeux8ba1}}

\begin{longtable}[]{@{}ll@{}}
\toprule
索引 & 作用 \\
\midrule
\endhead
\texttt{idx\_movies\_director\_id} & 支持导演详情页快速查找导演作品 \\
\texttt{idx\_movies\_release\_year} & 支持时间线和年代统计 \\
\texttt{idx\_movies\_rating} & 支持高分排序和推荐 \\
\texttt{idx\_movies\_status\_year} & 支持按状态和年份筛选 \\
\texttt{idx\_movie\_genres\_genre\_id} & 支持类型路由和类型筛选 \\
\texttt{idx\_movie\_people\_movie\_role} & 支持详情页演员表排序 \\
\texttt{idx\_user\_movie\_states\_user\_flags} &
支持收藏中心和已看统计 \\
\texttt{idx\_reviews\_movie\_id} & 支持电影详情和影评墙按时间读取影评 \\
\texttt{idx\_watch\_plans\_user\_status} &
支持待看片单按状态和日期组织 \\
\texttt{idx\_watch\_logs\_user\_watched\_at} & 支持观影日历和日志倒序 \\
\texttt{idx\_recent\_history\_user\_viewed\_at} & 支持首页最近浏览 \\
\texttt{idx\_audit\_logs\_entity} & 支持审计记录追踪 \\
\bottomrule
\end{longtable}

\begin{center}\rule{0.5\linewidth}{0.5pt}\end{center}

\hypertarget{api-ux8bbeux8ba1}{%
\section{API 设计}\label{api-ux8bbeux8ba1}}

\hypertarget{ux63a5ux53e3ux603bux89c8}{%
\subsection{接口总览}\label{ux63a5ux53e3ux603bux89c8}}

\begin{longtable}[]{@{}llll@{}}
\toprule
URL & 方法 & 功能 & 是否需要登录 \\
\midrule
\endhead
\texttt{/api/health} & GET & API 健康检查 & 否 \\
\texttt{/api/movies} & GET & 获取电影列表，支持标题和类型查询 & 否 \\
\texttt{/api/movies/:id} & GET & 获取电影详情 & 否 \\
\texttt{/api/movies} & POST & 新增电影 & 是 \\
\texttt{/api/movies/:id} & PUT & 更新电影 & 是 \\
\texttt{/api/movies/:id} & DELETE & 删除电影 & 是 \\
\texttt{/api/directors} & GET & 获取导演列表 & 否 \\
\texttt{/api/directors/:id} & GET & 获取导演详情 & 否 \\
\texttt{/api/directors/:id/movies} & GET & 获取导演作品 & 否 \\
\texttt{/api/directors} & POST & 新增导演 & 是 \\
\texttt{/api/directors/:id} & DELETE & 删除导演 & 是 \\
\bottomrule
\end{longtable}

\hypertarget{ux7535ux5f71ux5217ux8868ux63a5ux53e3}{%
\subsection{电影列表接口}\label{ux7535ux5f71ux5217ux8868ux63a5ux53e3}}

\begin{verbatim}
GET /api/movies?title=inception&genre=科幻
\end{verbatim}

接口职责：

\begin{itemize}
\tightlist
\item
  无参数时返回全部电影。
\item
  \texttt{title} 同时匹配片名和导演名。
\item
  \texttt{genre} 同时匹配主类型字段和多类型数组。
\item
  返回值是电影对象数组，前端会继续做规范化、媒体 URL
  检查和本地状态合并。
\end{itemize}

响应示例：

\begin{Shaded}
\begin{Highlighting}[]
\OtherTok{[}
  \FunctionTok{\{}
    \DataTypeTok{"id"}\FunctionTok{:} \DecValTok{2}\FunctionTok{,}
    \DataTypeTok{"title"}\FunctionTok{:} \StringTok{"盗梦空间 (Inception)"}\FunctionTok{,}
    \DataTypeTok{"director"}\FunctionTok{:} \StringTok{"克里斯托弗·诺兰"}\FunctionTok{,}
    \DataTypeTok{"genre"}\FunctionTok{:} \StringTok{"科幻"}\FunctionTok{,}
    \DataTypeTok{"genres"}\FunctionTok{:} \OtherTok{[}\StringTok{"科幻"}\OtherTok{,} \StringTok{"动作"}\OtherTok{,} \StringTok{"悬疑"}\OtherTok{]}\FunctionTok{,}
    \DataTypeTok{"rating"}\FunctionTok{:} \FloatTok{9.3}\FunctionTok{,}
    \DataTypeTok{"releaseYear"}\FunctionTok{:} \DecValTok{2010}\FunctionTok{,}
    \DataTypeTok{"posterUrl"}\FunctionTok{:} \StringTok{"https://image.tmdb.org/t/p/w500/89W962aAnPS3N3BdKgy2BvUhnCh.jpg"}
  \FunctionTok{\}}
\OtherTok{]}
\end{Highlighting}
\end{Shaded}

\hypertarget{ux7535ux5f71ux65b0ux589eux63a5ux53e3}{%
\subsection{电影新增接口}\label{ux7535ux5f71ux65b0ux589eux63a5ux53e3}}

\begin{verbatim}
POST /api/movies
Content-Type: application/json
\end{verbatim}

请求体示例：

\begin{Shaded}
\begin{Highlighting}[]
\FunctionTok{\{}
  \DataTypeTok{"title"}\FunctionTok{:} \StringTok{"沙丘2 (Dune: Part Two)"}\FunctionTok{,}
  \DataTypeTok{"director"}\FunctionTok{:} \StringTok{"丹尼斯·维伦纽瓦"}\FunctionTok{,}
  \DataTypeTok{"genre"}\FunctionTok{:} \StringTok{"科幻"}\FunctionTok{,}
  \DataTypeTok{"genres"}\FunctionTok{:} \OtherTok{[}\StringTok{"科幻"}\OtherTok{,} \StringTok{"冒险"}\OtherTok{,} \StringTok{"剧情"}\OtherTok{]}\FunctionTok{,}
  \DataTypeTok{"rating"}\FunctionTok{:} \FloatTok{8.5}\FunctionTok{,}
  \DataTypeTok{"releaseYear"}\FunctionTok{:} \DecValTok{2024}\FunctionTok{,}
  \DataTypeTok{"releaseDate"}\FunctionTok{:} \StringTok{"2024{-}03{-}01"}\FunctionTok{,}
  \DataTypeTok{"duration"}\FunctionTok{:} \DecValTok{166}\FunctionTok{,}
  \DataTypeTok{"posterUrl"}\FunctionTok{:} \StringTok{"https://image.tmdb.org/t/p/w500/9uaCR4HEZqxUqgORq0uZqTNm43G.jpg"}
\FunctionTok{\}}
\end{Highlighting}
\end{Shaded}

接口规则：

\begin{itemize}
\tightlist
\item
  \texttt{title} 必填。
\item
  后端生成唯一 \texttt{id}。
\item
  兼容单类型 \texttt{genre} 和多类型 \texttt{genres}。
\item
  \texttt{rating} 转为数值。
\item
  \texttt{releaseDate} 缺省时可由 \texttt{releaseYear} 补齐。
\end{itemize}

\hypertarget{ux5bfcux6f14ux63a5ux53e3}{%
\subsection{导演接口}\label{ux5bfcux6f14ux63a5ux53e3}}

导演接口提供导演库和导演详情的数据来源。\texttt{/api/directors/:id/movies}
用于导演详情页展示该导演作品集。前端 \texttt{DirectorService} 会在 API
不可用时从本地电影库按导演聚合，保证导演库仍可使用。

\hypertarget{ux6743ux9650ux7b56ux7565}{%
\subsection{权限策略}\label{ux6743ux9650ux7b56ux7565}}

系统权限边界采用``读公开、写受控''的设计：

\begin{itemize}
\tightlist
\item
  读接口：电影列表、电影详情、导演列表、导演详情、健康检查可公开访问。
\item
  写接口：新增、更新、删除电影和导演需要登录。
\item
  前端权限：\texttt{/add} 页面由 \texttt{authGuard}
  控制，未登录会跳回仪表盘并提示用户。
\item
  后端权限：正式数据库版本应在 POST、PUT、DELETE
  接口处校验用户身份和角色。
\item
  审计策略：关键写操作应记录到
  \texttt{audit\_logs}，保存操作人、实体、修改前和修改后的 JSON。
\end{itemize}

\begin{center}\rule{0.5\linewidth}{0.5pt}\end{center}

\hypertarget{ux524dux7aefux8bbeux8ba1}{%
\section{前端设计}\label{ux524dux7aefux8bbeux8ba1}}

\hypertarget{ux8defux7531ux7ed3ux6784}{%
\subsection{路由结构}\label{ux8defux7531ux7ed3ux6784}}

\begin{longtable}[]{@{}lll@{}}
\toprule
路径 & 页面 & 设计说明 \\
\midrule
\endhead
\texttt{/dashboard} & 仪表盘 & 系统首页，总览片库状态和快捷入口 \\
\texttt{/movies} & 电影库 & 搜索、筛选、排序、分页和视图切换 \\
\texttt{/movies/genre/:genre} & 分类浏览 &
参数化类型路由，复用电影库页面 \\
\texttt{/movies/:id/info} & 电影详情信息页 & 详情页默认子路由 \\
\texttt{/movies/:id/cast} & 演员表 & 详情页子路由，展示演职人员 \\
\texttt{/add} & 添加电影 & 受路由守卫保护 \\
\texttt{/directors} & 导演库 & 导演列表、搜索和作品入口 \\
\texttt{/directors/:id} & 导演详情 & 导演档案和作品聚合 \\
\texttt{/explore} & 探索影库 & 沉浸式浏览和快速扫片 \\
\texttt{/favorites} & 收藏中心 & 收藏电影与已看回看 \\
\texttt{/timeline} & 时间线 & 按年代组织电影 \\
\texttt{/recommendations} & 推荐页 & 基于类型和导演的推荐 \\
\texttt{/random} & 随机选片 & 快速抽取候选电影 \\
\texttt{/compare} & 电影对比 & 两部电影并列比较 \\
\texttt{/calendar} & 观影日历 & 按日期展示观看行为 \\
\texttt{/reviews} & 影评墙 & 管理用户评论和点赞 \\
\texttt{/watch-plans} & 待看片单 & 管理计划观看电影 \\
\texttt{/watch-logs} & 观影日志 & 记录观看场景和情绪 \\
\texttt{/smart-picks} & 智能选片 & 条件筛选和推荐预设 \\
\texttt{/director-atlas} & 导演图谱 & 导演维度统计 \\
\texttt{/mood-atlas} & 情绪图谱 & 情绪标签统计 \\
\texttt{/marathon} & 连看规划 & 按总时长组合片单 \\
\texttt{/taste-dna} & 偏好画像 & 汇总用户偏好 \\
\texttt{/scene-board} & 氛围策展 & 按视觉气质和主题组织电影 \\
\texttt{/archive-health} & 片库审计 & 检查资料完整度和 URL 质量 \\
\texttt{/about} & 关于页 & 系统概况、服务状态和数据管理入口 \\
\bottomrule
\end{longtable}

\hypertarget{ux7ec4ux4ef6ux7ed3ux6784}{%
\subsection{组件结构}\label{ux7ec4ux4ef6ux7ed3ux6784}}

\begin{verbatim}
src/app/
  pages/
    dashboard-page/
    movie-list-page/
    movie-detail-page/
    movie-add-page/
    director-list-page/
    director-detail-page/
    about-page/
  components/
    movie-dashboard/
    movie-list/
    movie-detail-info/
    movie-detail-cast/
    movie-form/
    movie-favorites/
    movie-timeline/
    movie-recommendations/
    movie-random/
    movie-compare/
    movie-calendar/
    movie-review-wall/
    watch-plans/
    watch-logs/
    smart-picks/
    director-atlas/
    mood-atlas/
    marathon-planner/
    taste-dna/
    scene-board/
    archive-health/
    command-palette/
    data-management/
\end{verbatim}

页面组件负责路由入口和页面容器，功能组件负责具体业务。比如电影详情页由父级
\texttt{movie-detail-page} 承载电影上下文，\texttt{movie-detail-info} 和
\texttt{movie-detail-cast} 作为子路由分别展示基本信息和演员表。

\hypertarget{ux670dux52a1ux5c42ux8bbeux8ba1}{%
\subsection{服务层设计}\label{ux670dux52a1ux5c42ux8bbeux8ba1}}

\begin{longtable}[]{@{}ll@{}}
\toprule
服务 & 职责 \\
\midrule
\endhead
\texttt{MovieService} & 电影主数据、搜索、收藏、已看、评分、HTTP
同步和本地降级 \\
\texttt{DirectorService} & 导演列表、导演详情、导演作品、API
与本地聚合 \\
\texttt{MovieStateService} & 聚合多个服务形成页面 ViewModel \\
\texttt{RecentHistoryService} & 记录最近浏览电影 \\
\texttt{ReviewStoreService} & 管理影评、点赞和影评列表 \\
\texttt{WatchPlanService} & 管理待看片单、状态和优先级 \\
\texttt{WatchLogService} & 管理观影日志、地点、评分和情绪标签 \\
\texttt{SmartPicksService} & 管理选片预设和智能推荐条件 \\
\texttt{DataPortService} & 导入导出电影、影评、日志、计划和预设 \\
\texttt{AuthService} & 管理演示登录状态 \\
\texttt{MessageService} & 用户可见消息 \\
\texttt{LoggerService} & 系统运行日志 \\
\bottomrule
\end{longtable}

\hypertarget{ux72b6ux6001ux6d41ux8bbeux8ba1}{%
\subsection{状态流设计}\label{ux72b6ux6001ux6d41ux8bbeux8ba1}}

\begin{center}
\centering
\includegraphics[width=0.96\linewidth,height=0.56\textheight,keepaspectratio]{output/pdf/diagrams/report-diagram-08.png}
\captionof{figure}{状态流设计}
\end{center}

状态流的核心是把分散的业务服务聚合为页面
ViewModel。页面组件不直接拼装复杂统计，而是订阅
\texttt{dashboardVm\$}、\texttt{movieListVm\$}、\texttt{movieDetailVm\$}
等 Observable，从而保持模板清晰。

\hypertarget{ux5a92ux4f53-url-ux8bbeux8ba1}{%
\subsection{媒体 URL 设计}\label{ux5a92ux4f53-url-ux8bbeux8ba1}}

电影海报、背景图和导演头像使用真实远程 URL。系统会对图片 URL
进行规范化：

\begin{itemize}
\tightlist
\item
  TMDB 图片统一转换为稳定尺寸路径。
\item
  Wikimedia 图片统一走可追溯公开图片地址。
\item
  旧的随机占位图或不稳定图源不会作为最终展示依据。
\item
  图片加载失败时使用通用电影素材兜底，而不是生成伪造电影图。
\item
  预告片使用``片名 + 年份 + official trailer''的 YouTube
  搜索链接，保证和电影标题相关。
\end{itemize}

\begin{center}\rule{0.5\linewidth}{0.5pt}\end{center}

\hypertarget{ux6838ux5fc3ux4ee3ux7801ux89e3ux6790}{%
\section{核心代码解析}\label{ux6838ux5fc3ux4ee3ux7801ux89e3ux6790}}

\hypertarget{ux8fd0ux884cux65f6-api-ux914dux7f6e}{%
\subsection{运行时 API 配置}\label{ux8fd0ux884cux65f6-api-ux914dux7f6e}}

代表文件：

\begin{itemize}
\tightlist
\item
  \texttt{src/app/config/api.config.ts}
\item
  \texttt{src/assets/runtime-config.js}
\end{itemize}

\begin{Shaded}
\begin{Highlighting}[]
\NormalTok{export const API\_CONFIG }\OperatorTok{=} \KeywordTok{new} \FunctionTok{InjectionToken}\OperatorTok{\textless{}}\NormalTok{ApiConfig}\OperatorTok{\textgreater{}}\NormalTok{(}\StringTok{\textquotesingle{}cinemaflow.api.config\textquotesingle{}}\OperatorTok{,}\NormalTok{ \{}
\NormalTok{  providedIn}\OperatorTok{:} \StringTok{\textquotesingle{}root\textquotesingle{}}\OperatorTok{,}
\NormalTok{  factory}\OperatorTok{:}\NormalTok{ () }\KeywordTok{=\textgreater{}}\NormalTok{ (\{}
\NormalTok{    baseUrl}\OperatorTok{:}\NormalTok{ window}\OperatorTok{.}\AttributeTok{\_\_CINEMAFLOW\_CONFIG\_\_}\OperatorTok{?.}\AttributeTok{apiBaseUrl} \OperatorTok{||} \StringTok{\textquotesingle{}/api\textquotesingle{}}
\NormalTok{  \})}
\NormalTok{\})}\OperatorTok{;}

\NormalTok{export }\KeywordTok{function} \FunctionTok{buildApiUrl}\NormalTok{(config}\OperatorTok{:}\NormalTok{ ApiConfig}\OperatorTok{,}\NormalTok{ path}\OperatorTok{:} \DataTypeTok{string}\NormalTok{)}\OperatorTok{:} \DataTypeTok{string}\NormalTok{ \{}
\NormalTok{  const baseUrl }\OperatorTok{=}\NormalTok{ config}\OperatorTok{.}\AttributeTok{baseUrl}\OperatorTok{.}\FunctionTok{replace}\NormalTok{(}\SpecialStringTok{/}\SpecialCharTok{\textbackslash{}/+$}\SpecialStringTok{/}\OperatorTok{,} \StringTok{\textquotesingle{}\textquotesingle{}}\NormalTok{)}\OperatorTok{;}
\NormalTok{  const normalizedPath }\OperatorTok{=}\NormalTok{ path}\OperatorTok{.}\FunctionTok{startsWith}\NormalTok{(}\StringTok{\textquotesingle{}/\textquotesingle{}}\NormalTok{) }\OperatorTok{?}\NormalTok{ path }\OperatorTok{:} \VerbatimStringTok{\textasciigrave{}/}\SpecialCharTok{$\{}\NormalTok{path}\SpecialCharTok{\}}\VerbatimStringTok{\textasciigrave{}}\OperatorTok{;}
  \FunctionTok{return} \VerbatimStringTok{\textasciigrave{}}\SpecialCharTok{$\{}\NormalTok{baseUrl}\SpecialCharTok{\}$\{}\NormalTok{normalizedPath}\SpecialCharTok{\}}\VerbatimStringTok{\textasciigrave{}}\OperatorTok{;}
\NormalTok{\}}
\end{Highlighting}
\end{Shaded}

代码说明：

\begin{itemize}
\tightlist
\item
  \texttt{InjectionToken} 让 API 配置成为 Angular 依赖注入的一部分。
\item
  \texttt{runtime-config.js} 在浏览器运行时提供
  \texttt{apiBaseUrl}，因此前端构建产物不绑定固定后端地址。
\item
  \texttt{buildApiUrl()}
  统一处理路径拼接，避免服务层散落字符串拼接逻辑。
\item
  \texttt{MovieService} 和 \texttt{DirectorService} 都通过该配置生成 API
  地址，服务层边界更清晰。
\end{itemize}

\hypertarget{movieservice-ux7684-http-ux540cux6b65ux4e0eux672cux5730ux964dux7ea7}{%
\subsection{MovieService 的 HTTP
同步与本地降级}\label{movieservice-ux7684-http-ux540cux6b65ux4e0eux672cux5730ux964dux7ea7}}

代表文件：\texttt{src/app/services/movie.service.ts}

\begin{Shaded}
\begin{Highlighting}[]
\NormalTok{private async }\FunctionTok{prefetchFromApiIfAvailable}\NormalTok{()}\OperatorTok{:} \BuiltInTok{Promise}\OperatorTok{\textless{}}\DataTypeTok{void}\OperatorTok{\textgreater{}}\NormalTok{ \{}
  \FunctionTok{if}\NormalTok{ (}\OperatorTok{!}\NormalTok{this}\OperatorTok{.}\AttributeTok{http}\NormalTok{) \{}
\NormalTok{    return}\OperatorTok{;}
\NormalTok{  \}}

\NormalTok{  this}\OperatorTok{.}\AttributeTok{http}\OperatorTok{.}\FunctionTok{get}\OperatorTok{\textless{}}\NormalTok{Movie[]}\OperatorTok{\textgreater{}}\NormalTok{(this}\OperatorTok{.}\AttributeTok{apiUrl}\NormalTok{)}\OperatorTok{.}\FunctionTok{pipe}\NormalTok{(}
    \FunctionTok{map}\NormalTok{(movies }\KeywordTok{=\textgreater{}}\NormalTok{ Array}\OperatorTok{.}\FunctionTok{isArray}\NormalTok{(movies) }\OperatorTok{?}\NormalTok{ movies}\OperatorTok{.}\FunctionTok{map}\NormalTok{(movie }\KeywordTok{=\textgreater{}}\NormalTok{ this}\OperatorTok{.}\FunctionTok{normalizeMovie}\NormalTok{(movie)) }\OperatorTok{:}\NormalTok{ [])}\OperatorTok{,}
    \FunctionTok{tap}\NormalTok{(movies }\KeywordTok{=\textgreater{}}\NormalTok{ \{}
      \FunctionTok{if}\NormalTok{ (movies}\OperatorTok{.}\AttributeTok{length} \OperatorTok{===} \DecValTok{0}\NormalTok{) \{}
\NormalTok{        return}\OperatorTok{;}
\NormalTok{      \}}

\NormalTok{      this}\OperatorTok{.}\FunctionTok{commitMovies}\NormalTok{(}
\NormalTok{        this}\OperatorTok{.}\FunctionTok{mergeApiMovies}\NormalTok{(movies)}\OperatorTok{,}
        \VerbatimStringTok{\textasciigrave{}Synced }\SpecialCharTok{$\{}\NormalTok{movies}\OperatorTok{.}\AttributeTok{length}\SpecialCharTok{\}}\VerbatimStringTok{ movies from Flask API\textasciigrave{}}
\NormalTok{      )}\OperatorTok{;}
\NormalTok{    \})}\OperatorTok{,}
    \FunctionTok{catchError}\NormalTok{(error }\KeywordTok{=\textgreater{}}\NormalTok{ \{}
\NormalTok{      this}\OperatorTok{.}\AttributeTok{logger}\OperatorTok{.}\FunctionTok{warn}\NormalTok{(}\VerbatimStringTok{\textasciigrave{}Flask movie API unavailable, using local catalog: }\SpecialCharTok{$\{}\NormalTok{this}\OperatorTok{.}\FunctionTok{describeHttpError}\NormalTok{(error)}\SpecialCharTok{\}}\VerbatimStringTok{\textasciigrave{}}\NormalTok{)}\OperatorTok{;}
\NormalTok{      return }\FunctionTok{of}\NormalTok{([])}\OperatorTok{;}
\NormalTok{    \})}
\NormalTok{  )}\OperatorTok{.}\FunctionTok{subscribe}\NormalTok{()}\OperatorTok{;}
\NormalTok{\}}
\end{Highlighting}
\end{Shaded}

代码说明：

\begin{itemize}
\tightlist
\item
  服务初始化后会尝试从 Flask API 获取电影数据。
\item
  API 返回的数据先进入 \texttt{normalizeMovie()}，统一日期、类型、导演
  ID、媒体 URL 和预告片 URL。
\item
  \texttt{mergeApiMovies()} 按``片名 + 年份''合并，避免重复电影。
\item
  后端不可用时不阻塞页面，系统继续使用本地片库，保证演示和浏览可用。
\item
  日志进入 \texttt{LoggerService}，关于页面可以展示运行状态。
\end{itemize}

\hypertarget{rxjs-ux9875ux9762-viewmodel}{%
\subsection{RxJS 页面 ViewModel}\label{rxjs-ux9875ux9762-viewmodel}}

代表文件：\texttt{src/app/services/movie-state.service.ts}

\begin{Shaded}
\begin{Highlighting}[]
\FunctionTok{movieListVm$}\NormalTok{(queryParamMap$}\OperatorTok{:}\NormalTok{ Observable}\OperatorTok{\textless{}}\NormalTok{ParamMap}\OperatorTok{\textgreater{},}\NormalTok{ routeParamMap$}\OperatorTok{?:}\NormalTok{ Observable}\OperatorTok{\textless{}}\NormalTok{ParamMap}\OperatorTok{\textgreater{}}\NormalTok{)}\OperatorTok{:}\NormalTok{ Observable}\OperatorTok{\textless{}}\NormalTok{MovieListViewModel}\OperatorTok{\textgreater{}}\NormalTok{ \{}
\NormalTok{  return }\FunctionTok{combineLatest}\NormalTok{([}
\NormalTok{    this}\OperatorTok{.}\AttributeTok{movieService}\OperatorTok{.}\AttributeTok{movies$}\OperatorTok{,}
\NormalTok{    queryParamMap$}\OperatorTok{,}
\NormalTok{    routeParamMap$ }\OperatorTok{??}\NormalTok{ queryParamMap$}
\NormalTok{  ])}\OperatorTok{.}\FunctionTok{pipe}\NormalTok{(}
    \FunctionTok{map}\NormalTok{(([movies}\OperatorTok{,}\NormalTok{ queryParamMap}\OperatorTok{,}\NormalTok{ routeParamMap]) }\KeywordTok{=\textgreater{}}\NormalTok{ \{}
\NormalTok{      const genres }\OperatorTok{=} \FunctionTok{collectMovieGenres}\NormalTok{(movies)}\OperatorTok{;}
\NormalTok{      const routeGenre }\OperatorTok{=}\NormalTok{ routeParamMap}\OperatorTok{.}\FunctionTok{get}\NormalTok{(}\StringTok{\textquotesingle{}genre\textquotesingle{}}\NormalTok{)}\OperatorTok{?.}\FunctionTok{trim}\NormalTok{() }\OperatorTok{??} \StringTok{\textquotesingle{}\textquotesingle{}}\OperatorTok{;}
\NormalTok{      const queryState }\OperatorTok{=} \FunctionTok{parseMovieQueryState}\NormalTok{(queryParamMap}\OperatorTok{,}\NormalTok{ genres)}\OperatorTok{;}
\NormalTok{      const nextQueryState }\OperatorTok{=}\NormalTok{ routeGenre }\OperatorTok{\&\&}\NormalTok{ genres}\OperatorTok{.}\FunctionTok{includes}\NormalTok{(routeGenre)}
        \OperatorTok{?}\NormalTok{ \{ }\OperatorTok{...}\NormalTok{queryState}\OperatorTok{,}\NormalTok{ genre}\OperatorTok{:}\NormalTok{ routeGenre \}}
        \OperatorTok{:}\NormalTok{ queryState}\OperatorTok{;}
\NormalTok{      const filteredMovies }\OperatorTok{=} \FunctionTok{filterMoviesByQueryState}\NormalTok{(movies}\OperatorTok{,}\NormalTok{ nextQueryState)}\OperatorTok{;}
\NormalTok{      const pagination }\OperatorTok{=} \FunctionTok{paginateItems}\NormalTok{(filteredMovies}\OperatorTok{,}\NormalTok{ nextQueryState}\OperatorTok{.}\AttributeTok{page}\OperatorTok{,}\NormalTok{ nextQueryState}\OperatorTok{.}\AttributeTok{pageSize}\NormalTok{)}\OperatorTok{;}

\NormalTok{      return \{}
\NormalTok{        movies}\OperatorTok{,}
\NormalTok{        filteredMovies}\OperatorTok{,}
\NormalTok{        visibleMovies}\OperatorTok{:}\NormalTok{ pagination}\OperatorTok{.}\AttributeTok{items}\OperatorTok{,}
\NormalTok{        genres}\OperatorTok{,}
\NormalTok{        queryState}\OperatorTok{:}\NormalTok{ nextQueryState}\OperatorTok{,}
\NormalTok{        summaryLabel}\OperatorTok{:} \VerbatimStringTok{\textasciigrave{}}\SpecialCharTok{$\{}\NormalTok{filteredMovies}\OperatorTok{.}\AttributeTok{length}\SpecialCharTok{\}}\VerbatimStringTok{ / }\SpecialCharTok{$\{}\NormalTok{movies}\OperatorTok{.}\AttributeTok{length}\SpecialCharTok{\}}\VerbatimStringTok{ 部\textasciigrave{}}\OperatorTok{,}
\NormalTok{        page}\OperatorTok{:}\NormalTok{ pagination}\OperatorTok{.}\AttributeTok{page}\OperatorTok{,}
\NormalTok{        pageSize}\OperatorTok{:}\NormalTok{ pagination}\OperatorTok{.}\AttributeTok{pageSize}\OperatorTok{,}
\NormalTok{        totalPages}\OperatorTok{:}\NormalTok{ pagination}\OperatorTok{.}\AttributeTok{totalPages}\OperatorTok{,}
\NormalTok{        startItem}\OperatorTok{:}\NormalTok{ pagination}\OperatorTok{.}\AttributeTok{startItem}\OperatorTok{,}
\NormalTok{        endItem}\OperatorTok{:}\NormalTok{ pagination}\OperatorTok{.}\AttributeTok{endItem}
\NormalTok{      \}}\OperatorTok{;}
\NormalTok{    \})}\OperatorTok{,}
    \FunctionTok{shareReplay}\NormalTok{(\{ bufferSize}\OperatorTok{:} \DecValTok{1}\OperatorTok{,}\NormalTok{ refCount}\OperatorTok{:}\NormalTok{ true \})}
\NormalTok{  )}\OperatorTok{;}
\NormalTok{\}}
\end{Highlighting}
\end{Shaded}

代码说明：

\begin{itemize}
\tightlist
\item
  \texttt{combineLatest()} 同时监听电影数据、URL 查询参数和路由参数。
\item
  \texttt{/movies} 和 \texttt{/movies/genre/:genre} 共用同一套列表
  ViewModel。
\item
  筛选、分页、类型解析都放在服务层，页面组件只负责展示和交互。
\item
  \texttt{shareReplay()} 让多个订阅复用同一份计算结果，减少重复计算。
\end{itemize}

\hypertarget{ux8defux7531ux5b88ux536b}{%
\subsection{路由守卫}\label{ux8defux7531ux5b88ux536b}}

代表文件：\texttt{src/app/guards/auth.guard.ts}

\begin{Shaded}
\begin{Highlighting}[]
\NormalTok{export const authGuard}\OperatorTok{:}\NormalTok{ CanActivateFn }\OperatorTok{=}\NormalTok{ () }\KeywordTok{=\textgreater{}}\NormalTok{ \{}
\NormalTok{  const authService }\OperatorTok{=} \FunctionTok{inject}\NormalTok{(AuthService)}\OperatorTok{;}
\NormalTok{  const messageService }\OperatorTok{=} \FunctionTok{inject}\NormalTok{(MessageService)}\OperatorTok{;}
\NormalTok{  const router }\OperatorTok{=} \FunctionTok{inject}\NormalTok{(Router)}\OperatorTok{;}

  \FunctionTok{if}\NormalTok{ (authService}\OperatorTok{.}\AttributeTok{isAuthenticated}\NormalTok{) \{}
\NormalTok{    return true}\OperatorTok{;}
\NormalTok{  \}}

\NormalTok{  messageService}\OperatorTok{.}\FunctionTok{warning}\NormalTok{(}\StringTok{\textquotesingle{}添加电影需要先登录。演示账号已固定为 admin / admin，可点击顶部按钮进入编辑模式。\textquotesingle{}}\OperatorTok{,} \StringTok{\textquotesingle{}路由守卫\textquotesingle{}}\NormalTok{)}\OperatorTok{;}
\NormalTok{  return router}\OperatorTok{.}\FunctionTok{createUrlTree}\NormalTok{([}\StringTok{\textquotesingle{}/dashboard\textquotesingle{}}\NormalTok{])}\OperatorTok{;}
\NormalTok{\}}\OperatorTok{;}
\end{Highlighting}
\end{Shaded}

代码说明：

\begin{itemize}
\tightlist
\item
  \texttt{CanActivateFn} 用于保护 \texttt{/add} 页面。
\item
  未登录时不显示空白页，而是通过消息服务提示用户。
\item
  守卫返回 \texttt{UrlTree}，由路由系统完成跳转。
\item
  正式后端版本应在写接口继续校验身份，前端守卫只负责用户体验和入口控制。
\end{itemize}

\begin{center}\rule{0.5\linewidth}{0.5pt}\end{center}

\hypertarget{ux9875ux9762ux622aux56feux4e0eux529fux80fdux8bf4ux660e}{%
\section{页面截图与功能说明}\label{ux9875ux9762ux622aux56feux4e0eux529fux80fdux8bf4ux660e}}

\begin{quote}
以下截图均为有效短视口截图，不使用超长全页图。截图文件位于
\texttt{docs/screenshots/}。
\end{quote}

\hypertarget{ux4eeaux8868ux76d8}{%
\subsection{仪表盘}\label{ux4eeaux8868ux76d8}}

\begin{center}
\centering
\includegraphics[width=0.94\linewidth,height=0.34\textheight,keepaspectratio]{docs/screenshots/dashboard.png}
\captionof{figure}{仪表盘}
\end{center}

仪表盘是系统入口，用于快速理解当前片库状态。页面展示总影片数、已观影数量、收藏数量、平均评分、类型分布、评分分布、导演排行、最近浏览和最近添加。用户进入系统后无需先搜索，就能知道片库整体质量、最近访问对象和下一步常用入口。

\hypertarget{ux7535ux5f71ux5e93}{%
\subsection{电影库}\label{ux7535ux5f71ux5e93}}

\begin{center}
\centering
\includegraphics[width=0.94\linewidth,height=0.34\textheight,keepaspectratio]{docs/screenshots/movies.png}
\captionof{figure}{电影库}
\end{center}

电影库是系统最核心的资料管理页面。它支持关键词搜索、类型筛选、排序、分页、表格/网格视图切换、已看和收藏状态识别。电影库承载``找到电影''的主流程，也为详情页、对比页、收藏中心和推荐页提供入口。

\hypertarget{ux641cux7d22ux72b6ux6001}{%
\subsection{搜索状态}\label{ux641cux7d22ux72b6ux6001}}

\begin{center}
\centering
\includegraphics[width=0.94\linewidth,height=0.34\textheight,keepaspectratio]{docs/screenshots/movies-search-inception.png}
\captionof{figure}{搜索状态}
\end{center}

搜索状态展示了 URL 查询参数和页面状态同步的能力。用户搜索
\texttt{Inception}
后，列表只显示匹配电影。该状态可以被复制为链接，适合在汇报、分享或测试时复现同一查询结果。

\hypertarget{ux7c7bux578bux8defux7531}{%
\subsection{类型路由}\label{ux7c7bux578bux8defux7531}}

\begin{center}
\centering
\includegraphics[width=0.94\linewidth,height=0.34\textheight,keepaspectratio]{docs/screenshots/movies-genre-sci-fi.png}
\captionof{figure}{类型路由}
\end{center}

类型路由使用 \texttt{/movies/genre/:genre}
进入指定类型片库。它和普通电影库复用同一个页面组件，但通过路由参数自动注入类型过滤条件。这说明系统没有为每个类型重复创建页面，而是使用参数化路由保持结构简洁。

\hypertarget{ux7535ux5f71ux8be6ux60c5ux4fe1ux606fux9875}{%
\subsection{电影详情信息页}\label{ux7535ux5f71ux8be6ux60c5ux4fe1ux606fux9875}}

\begin{center}
\centering
\includegraphics[width=0.94\linewidth,height=0.34\textheight,keepaspectratio]{docs/screenshots/movie-detail-info.png}
\captionof{figure}{电影详情信息页}
\end{center}

详情信息页展示单部电影的主视觉、海报、评分、类型、导演、简介、主演、个人评分、观影笔记、相邻电影和相似推荐。它是电影实体最完整的展示页面，也是收藏、已看、评分和日志等行为的上下文来源。

\hypertarget{ux7535ux5f71ux6f14ux5458ux8868}{%
\subsection{电影演员表}\label{ux7535ux5f71ux6f14ux5458ux8868}}

\begin{center}
\centering
\includegraphics[width=0.94\linewidth,height=0.34\textheight,keepaspectratio]{docs/screenshots/movie-detail-cast.png}
\captionof{figure}{电影演员表}
\end{center}

演员表是电影详情的子路由页面。父级详情页负责加载电影上下文，子页面只展示演员列表和角色信息。该设计让详情页可以继续扩展``剧照、幕后、奖项、影评''等更多子页面。

\hypertarget{ux6dfbux52a0ux7535ux5f71}{%
\subsection{添加电影}\label{ux6dfbux52a0ux7535ux5f71}}

\begin{center}
\centering
\includegraphics[width=0.94\linewidth,height=0.34\textheight,keepaspectratio]{docs/screenshots/add.png}
\captionof{figure}{添加电影}
\end{center}

添加电影页面用于录入新电影。表单覆盖片名、导演、上映日期、评分、类型、语言、时长、简介、演员、海报
URL、背景 URL 和预告片
URL。页面受登录守卫保护，避免普通浏览状态下误进入编辑入口。

\hypertarget{ux5bfcux6f14ux5e93}{%
\subsection{导演库}\label{ux5bfcux6f14ux5e93}}

\begin{center}
\centering
\includegraphics[width=0.94\linewidth,height=0.34\textheight,keepaspectratio]{docs/screenshots/directors.png}
\captionof{figure}{导演库}
\end{center}

导演库按导演聚合作品，展示头像、姓名、国籍、代表作品、活跃年代和风格标签。用户可以从导演维度进入作品集合，而不是只能从电影列表逐部查找。

\hypertarget{ux5bfcux6f14ux8be6ux60c5}{%
\subsection{导演详情}\label{ux5bfcux6f14ux8be6ux60c5}}

\begin{center}
\centering
\includegraphics[width=0.94\linewidth,height=0.34\textheight,keepaspectratio]{docs/screenshots/director-detail.png}
\captionof{figure}{导演详情}
\end{center}

导演详情页展示导演档案、风格描述、获奖信息、代表作品和作品列表。该页面强调``作者''和``作品''的关系，适合分析某位导演在片库中的风格和评分表现。

\hypertarget{ux63a2ux7d22ux5f71ux5e93}{%
\subsection{探索影库}\label{ux63a2ux7d22ux5f71ux5e93}}

\begin{center}
\centering
\includegraphics[width=0.94\linewidth,height=0.34\textheight,keepaspectratio]{docs/screenshots/explore.png}
\captionof{figure}{探索影库}
\end{center}

探索影库提供更沉浸的浏览方式，适合用户没有明确搜索目标时快速扫片。页面通过大图、评分和基础信息帮助用户快速判断是否进入详情。

\hypertarget{ux6536ux85cfux4e2dux5fc3}{%
\subsection{收藏中心}\label{ux6536ux85cfux4e2dux5fc3}}

\begin{center}
\centering
\includegraphics[width=0.94\linewidth,height=0.34\textheight,keepaspectratio]{docs/screenshots/favorites.png}
\captionof{figure}{收藏中心}
\end{center}

收藏中心汇总用户标记为收藏的电影。它相当于用户个人片库中的高价值子集，适合快速回看、推荐和偏好分析。收藏数据会被偏好画像和智能选片作为重要信号。

\hypertarget{ux65f6ux95f4ux7ebf}{%
\subsection{时间线}\label{ux65f6ux95f4ux7ebf}}

\begin{center}
\centering
\includegraphics[width=0.94\linewidth,height=0.34\textheight,keepaspectratio]{docs/screenshots/timeline.png}
\captionof{figure}{时间线}
\end{center}

时间线按上映年代组织电影，帮助用户观察片库覆盖的时间跨度。它能展示用户偏好的年代分布，也能发现某些年代的资料空白。

\hypertarget{ux63a8ux8350ux9875}{%
\subsection{推荐页}\label{ux63a8ux8350ux9875}}

\begin{center}
\centering
\includegraphics[width=0.94\linewidth,height=0.34\textheight,keepaspectratio]{docs/screenshots/recommendations.png}
\captionof{figure}{推荐页}
\end{center}

推荐页基于电影类型、导演和评分相近度生成推荐。它面向``看完一部后还想看类似电影''的场景，是从详情页延伸出的发现路径。

\hypertarget{ux968fux673aux9009ux7247}{%
\subsection{随机选片}\label{ux968fux673aux9009ux7247}}

\begin{center}
\centering
\includegraphics[width=0.94\linewidth,height=0.34\textheight,keepaspectratio]{docs/screenshots/random.png}
\captionof{figure}{随机选片}
\end{center}

随机选片用于解决临时决策问题。用户不想设置复杂条件时，可以直接让系统抽取候选电影。该功能适合轻量场景，也能和智能选片形成互补。

\hypertarget{ux7535ux5f71ux5bf9ux6bd4}{%
\subsection{电影对比}\label{ux7535ux5f71ux5bf9ux6bd4}}

\begin{center}
\centering
\includegraphics[width=0.94\linewidth,height=0.34\textheight,keepaspectratio]{docs/screenshots/compare.png}
\captionof{figure}{电影对比}
\end{center}

电影对比页面允许选择两部电影并列查看，从评分、时长、类型、导演、语言等维度比较。它适合用户在两部候选电影之间犹豫时做决策。

\hypertarget{ux89c2ux5f71ux65e5ux5386}{%
\subsection{观影日历}\label{ux89c2ux5f71ux65e5ux5386}}

\begin{center}
\centering
\includegraphics[width=0.94\linewidth,height=0.34\textheight,keepaspectratio]{docs/screenshots/calendar.png}
\captionof{figure}{观影日历}
\end{center}

观影日历基于观影日志，而不是电影上映日期。它展示用户真实观看行为在时间上的分布，可以看到哪天观看了什么、年度观看热度如何。

\hypertarget{ux5f71ux8bc4ux5899}{%
\subsection{影评墙}\label{ux5f71ux8bc4ux5899}}

\begin{center}
\centering
\includegraphics[width=0.94\linewidth,height=0.34\textheight,keepaspectratio]{docs/screenshots/reviews.png}
\captionof{figure}{影评墙}
\end{center}

影评墙用于沉淀文字评价。用户可以为电影写下作者、评分和评论内容，也可以浏览已有评论和点赞。影评是偏好画像和电影评价的重要补充。

\hypertarget{ux5f85ux770bux7247ux5355}{%
\subsection{待看片单}\label{ux5f85ux770bux7247ux5355}}

\begin{center}
\centering
\includegraphics[width=0.94\linewidth,height=0.34\textheight,keepaspectratio]{docs/screenshots/watch-plans.png}
\captionof{figure}{待看片单}
\end{center}

待看片单管理``未来想看''的电影。每条计划包含优先级、计划日期、观看场景、状态和备注。它把电影从普通收藏进一步推进到具体观看安排。

\hypertarget{ux89c2ux5f71ux65e5ux5fd7}{%
\subsection{观影日志}\label{ux89c2ux5f71ux65e5ux5fd7}}

\begin{center}
\centering
\includegraphics[width=0.94\linewidth,height=0.34\textheight,keepaspectratio]{docs/screenshots/watch-logs.png}
\captionof{figure}{观影日志}
\end{center}

观影日志记录实际观看事件，包括观看时间、地点、陪伴对象、会话评分、情绪标签和笔记。日志是观影日历、情绪图谱和偏好画像的基础数据。

\hypertarget{ux667aux80fdux9009ux7247}{%
\subsection{智能选片}\label{ux667aux80fdux9009ux7247}}

\begin{center}
\centering
\includegraphics[width=0.94\linewidth,height=0.34\textheight,keepaspectratio]{docs/screenshots/smart-picks.png}
\captionof{figure}{智能选片}
\end{center}

智能选片按时长、最低评分、类型、语言、是否包含已看、是否偏好收藏等条件生成候选电影。用户可以保存常用预设，例如``下班后轻松看''``周末长片''``只看高分科幻''。

\hypertarget{ux5bfcux6f14ux56feux8c31}{%
\subsection{导演图谱}\label{ux5bfcux6f14ux56feux8c31}}

\begin{center}
\centering
\includegraphics[width=0.94\linewidth,height=0.34\textheight,keepaspectratio]{docs/screenshots/director-atlas.png}
\captionof{figure}{导演图谱}
\end{center}

导演图谱按导演聚合作品数量、平均评分、收藏数、已看数和总时长。它把电影列表转换为导演维度的统计视图，适合观察片库中哪些导演占比更高。

\hypertarget{ux60c5ux7eeaux56feux8c31}{%
\subsection{情绪图谱}\label{ux60c5ux7eeaux56feux8c31}}

\begin{center}
\centering
\includegraphics[width=0.94\linewidth,height=0.34\textheight,keepaspectratio]{docs/screenshots/mood-atlas.png}
\captionof{figure}{情绪图谱}
\end{center}

情绪图谱读取观影日志里的情绪标签，统计治愈、热血、沉重、浪漫等标签分布。它帮助用户从``看了什么''进一步分析``为什么喜欢''和``观看时的情绪需求''。

\hypertarget{ux8fdeux770bux89c4ux5212}{%
\subsection{连看规划}\label{ux8fdeux770bux89c4ux5212}}

\begin{center}
\centering
\includegraphics[width=0.94\linewidth,height=0.34\textheight,keepaspectratio]{docs/screenshots/marathon.png}
\captionof{figure}{连看规划}
\end{center}

连看规划根据用户设置的总时长预算和类型偏好组合连续观影片单。它适合周末、假期或主题观影夜，解决``多部电影如何组合''的问题。

\hypertarget{ux504fux597dux753bux50cf}{%
\subsection{偏好画像}\label{ux504fux597dux753bux50cf}}

\begin{center}
\centering
\includegraphics[width=0.94\linewidth,height=0.34\textheight,keepaspectratio]{docs/screenshots/taste-dna.png}
\captionof{figure}{偏好画像}
\end{center}

偏好画像综合已看、收藏、个人评分、影评和观影日志，展示类型、导演、语言、年代等偏好。它是系统从资料管理走向个人分析的核心页面。

\hypertarget{ux6c1bux56f4ux7b56ux5c55}{%
\subsection{氛围策展}\label{ux6c1bux56f4ux7b56ux5c55}}

\begin{center}
\centering
\includegraphics[width=0.94\linewidth,height=0.34\textheight,keepaspectratio]{docs/screenshots/scene-board.png}
\captionof{figure}{氛围策展}
\end{center}

氛围策展按视觉气质、主题和情绪重新组织电影。它不只按类型分类，而是从``适合什么心情、什么夜晚、什么场景''出发形成主题片单。

\hypertarget{ux7247ux5e93ux5ba1ux8ba1}{%
\subsection{片库审计}\label{ux7247ux5e93ux5ba1ux8ba1}}

\begin{center}
\centering
\includegraphics[width=0.94\linewidth,height=0.34\textheight,keepaspectratio]{docs/screenshots/archive-health.png}
\captionof{figure}{片库审计}
\end{center}

片库审计用于检查电影资料完整度，包含海报 URL、背景
URL、简介、演员表、语言、片长等字段。它帮助资料维护者发现缺失项和低质量媒体资源。

\hypertarget{ux5173ux4e8eux9875ux9762}{%
\subsection{关于页面}\label{ux5173ux4e8eux9875ux9762}}

\begin{center}
\centering
\includegraphics[width=0.94\linewidth,height=0.34\textheight,keepaspectratio]{docs/screenshots/about.png}
\captionof{figure}{关于页面}
\end{center}

关于页展示系统技术栈、模块总览、服务状态、最近消息和运行日志。它既是系统说明页，也是观察服务状态和数据管理入口的集中页面。

\hypertarget{ux547dux4ee4ux9762ux677f}{%
\subsection{命令面板}\label{ux547dux4ee4ux9762ux677f}}

\begin{center}
\centering
\includegraphics[width=0.94\linewidth,height=0.34\textheight,keepaspectratio]{docs/screenshots/command-palette.png}
\captionof{figure}{命令面板}
\end{center}

命令面板提供全局快速跳转和搜索。用户可以搜索页面、电影或导演，并直接进入目标位置。它减少了多页面系统中的导航成本。

\hypertarget{ux6700ux8fd1ux6d4fux89c8ux533aux57df}{%
\subsection{最近浏览区域}\label{ux6700ux8fd1ux6d4fux89c8ux533aux57df}}

\begin{center}
\centering
\includegraphics[width=0.94\linewidth,height=0.34\textheight,keepaspectratio]{docs/screenshots/recent-history.png}
\captionof{figure}{最近浏览区域}
\end{center}

最近浏览区域展示用户近期打开过的电影详情。它帮助用户回到刚才看过的电影，尤其适合在对比、推荐和搜索之间频繁跳转时恢复上下文。

\hypertarget{ux6570ux636eux7ba1ux7406ux533aux57df}{%
\subsection{数据管理区域}\label{ux6570ux636eux7ba1ux7406ux533aux57df}}

\begin{center}
\centering
\includegraphics[width=0.94\linewidth,height=0.34\textheight,keepaspectratio]{docs/screenshots/data-management.png}
\captionof{figure}{数据管理区域}
\end{center}

数据管理区域支持导入导出片库、影评、最近浏览、待看片单、观影日志和智能预设。它为个人资料迁移和备份提供入口，也让系统具备离线演示和恢复能力。

\begin{center}\rule{0.5\linewidth}{0.5pt}\end{center}

\hypertarget{ux7cfbux7edfux95eeux9898ux4e0eux603bux7ed3}{%
\section{系统问题与总结}\label{ux7cfbux7edfux95eeux9898ux4e0eux603bux7ed3}}

\hypertarget{ux5173ux952eux8bbeux8ba1ux53d6ux820d}{%
\subsection{关键设计取舍}\label{ux5173ux952eux8bbeux8ba1ux53d6ux820d}}

\begin{longtable}[]{@{}ll@{}}
\toprule
设计点 & 取舍说明 \\
\midrule
\endhead
前端本地状态与后端 API 并存 &
保证浏览体验稳定，同时保留向数据库持久化演进的接口边界 \\
图片 URL 真实化 &
避免随机占位图造成资料不可信，优先使用可追溯远程资源 \\
电影状态拆为用户状态表 &
多用户场景下，同一部电影的收藏、已看和评分应因人而异 \\
页面 ViewModel 服务化 &
避免页面模板直接拼装复杂统计，让组件更专注展示 \\
详情页使用子路由 & 让基本信息、演员表等内容可以独立扩展 \\
片库审计独立成页 & 把资料质量变成可见功能，而不是隐藏在代码逻辑中 \\
\bottomrule
\end{longtable}

\hypertarget{ux5f53ux524dux7cfbux7edfux9002ux5408ux7684ux573aux666f}{%
\subsection{当前系统适合的场景}\label{ux5f53ux524dux7cfbux7edfux9002ux5408ux7684ux573aux666f}}

\begin{itemize}
\tightlist
\item
  个人电影资料管理。
\item
  电影课程项目展示。
\item
  前后端分离架构演示。
\item
  Angular 状态流和组件化设计展示。
\item
  关系型数据库建模说明。
\item
  影迷偏好分析和观影记录沉淀。
\end{itemize}

\hypertarget{ux540eux7eedux53efux6269ux5c55ux65b9ux5411}{%
\subsection{后续可扩展方向}\label{ux540eux7eedux53efux6269ux5c55ux65b9ux5411}}

\begin{itemize}
\tightlist
\item
  后端接入 SQLAlchemy，将 JSON 持久层替换为关系型数据库。
\item
  增加登录接口、JWT 或 Session，并在写接口做权限校验。
\item
  增加 OpenAPI 文档，提升接口可读性。
\item
  增加数据库迁移工具，例如 Alembic。
\item
  将影评、日志、计划、智能预设全部迁移到后端持久化。
\item
  增加端到端测试覆盖搜索、添加电影、收藏、日志、导入导出等路径。
\item
  为媒体 URL 增加定时健康检查和自动修复建议。
\end{itemize}

\hypertarget{ux603bux7ed3}{%
\subsection{总结}\label{ux603bux7ed3}}

CinemaFlow
的核心价值在于把电影资料、用户观影行为和偏好分析连接起来。它不是单纯的电影列表，而是一个完整的私人电影资料馆系统。系统通过前后端分离架构保持页面交互和后端数据边界清晰，通过关系型数据库设计表达电影、导演、演员、用户行为和审计之间的关系，通过多页面功能覆盖浏览、记录、推荐、分析和维护的完整链路。

\end{document}
